\chapter{Тетраэдрическая матричная кривизна и развилка калибровочного кадра}

Предыдущая глава отождествила полную дефектную связность с единым
семейным расслоением и оставила один вопрос: выводится ли тетраэдрический
конденсат из того же родителя. Перед настоящим гейтом были отдельно
проверены проектный журнал и первичная литература. Это позволило заменить
общий поиск потенциала одной точной матричной кривизной.

\section{Литературная и проектная граница}

В общей перенормируемой теории поля спина три $SO(3)$-симметрия допускает
два независимых квартичных инварианта. Поэтому сама симметрия не выбирает
тетраэдрическую фазу; требуется определённая область относительных
коэффициентов. Это показано в \path{arXiv:0910.4392} и в систематическом
анализе JHEP 08 (2017) 110.

Работа \path{arXiv:2607.12366} даёт более подходящую для проекта форму:
потенциал записывается как сумма неотрицательного тетраэдрического
инварианта и радиального квадрата. Ниже показано, что оба слагаемых можно
объединить в одну норму симметричной матричной кривизны.

Ранний проектный квадрат
\begin{equation}
 Q(H)=H^2-\frac1{\sqrt3}H-\frac14I_4
\end{equation}
действительно выбирал четыре примитивных проектора. Но ровно четыре
решения получались после ограничения $H$ диагональным бесследовым меню.
Поэтому этот результат нельзя без дополнительного доказательства считать
полным $SO(3)$-инвариантным механизмом спонтанного нарушения.

\section{Одна кривизна на полном представлении спина три}

Пусть
\begin{equation}
 T\in\operatorname{Sym}^3_0(E_{\rm fam})
\end{equation}
--- вещественное семимерное поле спина три. Определим
\begin{equation}
 A_{ij}(T)=T_{ikl}T_{jkl},
 \qquad
 \mu_T(T)=A(T)-\frac{v_T^2}{3}I_3.
 \label{eq:v6-tetrahedral-matrix-curvature}
\end{equation}
Если
\begin{equation}
 J_2=\Tr A,
 \qquad
 A_0=A-\frac{J_2}{3}I_3,
\end{equation}
то выполнено точное тождество
\begin{equation}
 \boxed{
 \Tr\mu_T^2
 =\Tr A_0^2+\frac13(J_2-v_T^2)^2.}
 \label{eq:v6-tetrahedral-curvature-decomposition}
\end{equation}
Случайный аудит на ста тензорах дал максимальный остаток
$7.2\cdot10^{-15}$.

Бесследовая часть $A_0$ является спин-две проекцией квадрата $T\otimes T$.
Литературный тетраэдрический инвариант пропорционален
\begin{equation}
 3\Tr A_0^2=3J_4-J_2^2.
\end{equation}
Центральная часть~\eqref{eq:v6-tetrahedral-curvature-decomposition}
фиксирует ненулевую амплитуду. Таким образом, один невзвешенный квадрат на
модуле
\begin{equation}
 \operatorname{Sym}_3(\mathbb R)simeq\mathbf1\oplus\mathbf5
\end{equation}
заменяет два независимо записанных квартичных инварианта.

Именно этот модуль уже появлялся в гейтах отображения момента Томов IV--V.
Прежний запрет идеала вырожденных форм сохраняется для конкретной
18-мерной KO6-цепи и её обычного дифференциального исчисления. Он не
является общим запретом на составную кривизну поля спина три, но требует
предъявить для неё новый допустимый родительский блок.

\section{Независимое поле спина три}

Возьмём четыре единичные тетраэдрические оси $n_a$ и канонический вакуум
\begin{equation}
 T_\star=\sum_{a=1}^4n_a^{\otimes3},
 \qquad
 v_T^2=\|T_\star\|^2=\frac{32}{9}.
 \label{eq:v6-tetrahedral-spin-three-vacuum}
\end{equation}
Для него
\begin{equation}
 A(T_\star)=\frac{v_T^2}{3}I_3,
 \qquad
 \mu_T(T_\star)=0.
\end{equation}

Полный гессиан функционала
\begin{equation}
 V_T=\tau_3(\mu_T^2)
 \label{eq:v6-independent-tetrahedral-potential}
\end{equation}
вычислен на всех семи компонентах $\operatorname{Sym}^3_0(\mathbb R^3)$.
Его спектр равен
\begin{equation}
 \left\{
 0^{(3)},
 \left(\frac{512}{405}\right)^{(3)},
 \frac{256}{81}
 \right\}.
 \label{eq:v6-independent-tetrahedral-hessian}
\end{equation}
Три нуля в точности совпадают с касательными к орбите $SO(3)/A_4$;
остаток действия гессиана на них равен $6.6\cdot10^{-15}$. Остальные
четыре направления строго положительны.

Двенадцать собственных вращений тетраэдра сохраняют $T_\star$ с остатком
не более $1.4\cdot10^{-15}$, поэтому собственный стабилизатор в $SO(3)$
равен $A_4$. Кинетический член $\|D_BT\|^2$ даёт трём компонентам
семейной связности одинаковые положительные массовые коэффициенты
$128/9$ в используемой нормировке.

Нормировка потенциального квадрата также не требует нового следа:
\begin{equation}
 \frac1{45}\Tr_{45}
 \left((\mu_T^*\mu_T)\otimes I_{15}\right)
 =\frac13\Tr_3(\mu_T^*\mu_T).
\end{equation}
Машинный остаток равен нулю.

\begin{theorem}[Локальный тетраэдрический проход]
Независимое поле $T\in\operatorname{Sym}^3_0(E_{\rm fam})$ с одним
квадратом $\tau_3(\mu_T^2)$ имеет вакуумную орбиту $SO(3)/A_4$, ровно три
калибровочных нуля и четыре положительные нормальные моды. Та же семейная
связность получает массу во всех трёх направлениях, а потенциальная норма
совпадает с ограничением общего физического следа.
\end{theorem}

Но семимерное поле ещё не выведено как независимая динамическая переменная
текущего конечного родителя. Математический проход поэтому не является
автоматическим физическим принятием нового поля.

\section{Составной четырёхосный кадр}

Проверим вариант без нового фундаментального поля. Пусть $X$ есть
$3\times4$-матрица, столбцами которой в вакууме служат оси $n_a$:
\begin{equation}
 X_\star=(n_1,n_2,n_3,n_4),
 \qquad
 T(X)=\sum_{a=1}^4x_a^{\otimes3}.
 \label{eq:v6-composite-tetrahedral-frame}
\end{equation}
Тогда $T(X_\star)=T_\star$. Кроме того,
\begin{equation}
 X_\star X_\star^T=\frac43I_3,
 \qquad
 X_\star^TX_\star=\frac43P_3,
\end{equation}
где $P_3=I_4-uu^T$ есть канонический проектор, ортогональный аффинному
синглету $u=(1,1,1,1)^T/2$.

Поэтому рассмотрим одну блочную кривизну
\begin{equation}
 \mathcal F_X=operatorname{diag}\left(
 XX^T-\frac43I_3,
 X^TX-\frac43P_3
 \right)
 \label{eq:v6-composite-frame-curvature}
\end{equation}
и её невзвешенную следовую норму. Никакая отдельная ось здесь не задана:
$P_3$ инвариантен относительно всего перестановочного действия $S_4$.

Полный двенадцатимерный гессиан $\Tr\mathcal F_X^2$ имеет спектр
\begin{equation}
 \left\{
 0^{(3)},
 \left(\frac{16}{3}\right)^{(3)},
 \left(\frac{64}{3}\right)^{(6)}
 \right\}.
 \label{eq:v6-composite-frame-hessian}
\end{equation}
Следовательно, девять нормальных направлений положительны, а три нуля
являются общими левыми вращениями кадра. Добавление составной кривизны
$\mu_T(T(X))$ не меняет нульность и поднимает только уже положительные
направления.

Этот вариант экономнее по полям, но изменяет физический тип остаточной
симметрии. Стабилизатор $X_\star$ относительно одного левого
калибровочного $SO(3)$ тривиален: из $gX_\star=X_\star$ следует $g=I$.
Двенадцать элементов $A_4$ возникают только как совместные пары
\begin{equation}
 R_\sigma X_\star=X_\star U_\sigma,
 \qquad \sigma\in A_4,
 \label{eq:v6-diagonal-tetrahedral-frame}
\end{equation}
где $U_\sigma$ переставляет четыре правые оси. Максимальный остаток
равен $3.9\cdot10^{-16}$.

\begin{proposition}[Составной проход только для диагонального кадра]
Четырёхосный частично-изометрический кадр выводит тетраэдрический тензор
без нового фундаментального поля и имеет положительный нормальный
гессиан. Но он не оставляет $A_4$ как подгруппу одного левого
калибровочного $SO(3)$. Группа $A_4$ сохраняется только как диагональная
калибровочно-кадровая или гранично фиксированная симметрия.
\end{proposition}

Это в точности возвращает старую развилку между истинно калибровочной и
гранично фиксированной семейной осью. Она больше не скрыта в коэффициентах:
оба варианта вычислены, но описывают разные физические архитектуры.

\section{Что закрыто и что осталось}

Получен сильный положительный результат: тетраэдрический конденсат не
требует произвольного многопараметрического потенциала. Полное
семимерное поле допускает один матричный квадрат с правильной орбитой и
положительным гессианом. Составной кадр даёт ещё более экономное
вариационное происхождение.

Однако единственного физического ответа пока нет:
\begin{enumerate}
 \item независимый $T$ даёт настоящую остаточную калибровочную $A_4$, но
 ещё не выведен как поле текущего родителя;
 \item составной $T(X)$ уже построен из существующего кадра, но оставляет
 только диагональную или гранично фиксированную $A_4$, а не остаточную
 подгруппу одного калибровочного $SO(3)$.
\end{enumerate}

\begin{nogocriterion}[Запрет смешения двух тетраэдрических чтений]
Нельзя использовать составной кадр как доказательство существования
остаточного калибровочного $A_4$ и нельзя использовать независимое поле
спина три как уже выведенную переменную конечного родителя. Следующий шаг
обязан выбрать и обосновать одну из этих физических ветвей.
\end{nogocriterion}

\section{Вердикт}

\begin{protocol}
\begin{enumerate}
 \item литературная матричная кривизна: \textbf{получена};
 \item один квадрат для формы и амплитуды: \textbf{получен};
 \item модуль кривизны: \textbf{$\mathbf1\oplus\mathbf5$};
 \item независимый семимерный гессиан: \textbf{$3$ нуля и $4$ плюса};
 \item стабилизатор независимого поля: \textbf{$A_4$};
 \item масса всех трёх семейных компонент: \textbf{положительна};
 \item общий след потенциального квадрата: \textbf{совпадает};
 \item независимое поле $T$ выведено текущим родителем: \textbf{нет};
 \item составной четырёхосный кадр: \textbf{проходит гессиан};
 \item новые фундаментальные компоненты в составной ветви: \textbf{нет};
 \item остаточная левая калибровочная группа составной ветви:
 \textbf{тривиальна};
 \item диагональная или гранично фиксированная $A_4$:
 \textbf{получена};
 \item единственная физическая ветвь: \textbf{не выбрана};
 \item рождение материи: \textbf{не закрыто};
 \item следующий гейт: \textbf{решение калибровочно-кадровой развилки}.
\end{enumerate}
\end{protocol}

Машинный сертификат сохранён в
\path{s2t_v6_bosonic_defect_tetrahedral_gauge_mass_parent_gate_results.json}.